\scanpage{17}
\mspage{4}
\deleted{(Le cas $I=\varnothing$ étant trivial).}

$1^\circ$) Cas $\alpha=0$ i.e. $I$ \deleted{réduit à un élément, \unsure{donc} $M(I^\circ)\overset{\mathrm{can}}{\simeq} M$. Alors \unsure{il est tout} immédiat que $C_I\unsure{\simeq|=}C$, $(TC)_I=TC$, $W_{I^\circ}=W$, et l'énoncé est trivial, tautologique.} \inserted{$=\varnothing$ : il est trivial}

$2^\circ$) Si c'est OK pour $I_\alpha$, aussi pour $I_{\alpha+1}$ \inserted{$=\{\beta\mid\beta\leq\alpha\}$}, \deleted{\unsure{Ici c'est le plus délicat, d'abord}} \inserted{$I_{\alpha+1}$ ayant un plus grand élément $\alpha$,} a) est contenu dans b) (pour $i=\alpha$, \unsure{à|=} plus \deleted{\unsure{grand}} élément de $I_{\alpha+1}$).

\deleted{Considérons alors le diagramme} \deleted{Alors c) résulte de l'hypothèse de récurrence, qui implique}
Pour prouver b), il reste à prouver que $u_\alpha\in C$. On regarde le diagramme $(\alpha)$ \inserted{pour $i_0=\alpha$.} Par \deleted{\unsure{ailleurs}} l'hyp. de récurrence \inserted{(pour a)} appliquée à $I_\alpha$, on a $\alpha_{\unsure{i_0|\alpha_0}}\in C$, donc par cochgt. de base $\beta_{i_0}\in C$, donc par composition $u_{i_0}=\widetilde{u}_{i_0}\beta_{i_0}\in C$.

Pour prouver c), il reste à prouver que $\varinjlim_{i<\alpha}u_i\in W$. \deleted{mais par hyp. de réc., \unsure{on} sait que les $u_i$ $(i<\alpha)$ sont dans $W$, \unsure{donc}} Mais par l'hyp. de réc. appliquée à $I_\alpha=\{i\mid i<\alpha\}$, on sait déjà que $(u_i)_{i<\alpha}\in (TC)_{I_\alpha}$, donc par la partie prouvée du corollaire ou de la prop., $\varinjlim_{i<\alpha}u_i\in TC$.

$3^\circ$) Si $\alpha$ un ordinal limite, et si a) b) c) vrais pour $I_\beta$ \unsure{avec} $\beta<\alpha$, aussi pour $I_\alpha$.

- Cette fois, b) est contenu dans l'hyp. de
% archive sheet number: 16

